(1) 二维 k 空间每单位面积态数密度为 \(A/(2\pi)^2\)，自旋简并为2。故填满半径 kF 的圆盘时
\[
N=2\frac A{(2\pi)^2}\pi k_F^2=\frac{Ak_F^2}{2\pi},\quad
k_F=\sqrt{2\pi n}.
\]
费米能 \(E_F=\hbar^2k_F^2/(2m^*)=\pi\hbar^2 n/m^*\)。单位实际面积、单位能量的态密度为
\[
D(E)=\frac1A\frac{dN(E)}{dE}
=\frac{m^*}{\pi\hbar^2},\quad E>0.
\]
常数性质来自二维累计态数 \(N(E)\propto k^2\propto E\)。

(2) 零温总内能可由能量态密度积分
\[
U=A\int_0^{E_F}E D(E)\,dE
=\frac{Am^*E_F^2}{2\pi\hbar^2}
=\frac12NE_F=\frac{\pi\hbar^2N^2}{2m^*A}.
\]
固定粒子数变化面积所得二维压强（单位为能量/面积）
\[
P_{\!2D}=-\left(\frac{\partial U}{\partial A}\right)_N
=\frac{\pi\hbar^2N^2}{2m^*A^2}=\frac12nE_F.
\]

(3) 化学势为 \(\mu=(\partial U/\partial N)_A=E_F\)，而二维体积模量对应的面积压缩模量为
\[
B_{\!2D}=-A(\partial P_{\!2D}/\partial A)_N
=2P_{\!2D}=nE_F.
\]
它是正值，体现占据动量态造成的 Pauli 压强而非显式电子间排斥势。

(4) 电子数密度和化学势关系 \(n=m^*\mu/(\pi\hbar^2)\)。因此
\[
\frac{C_q}{A}=e^2\frac{\partial n}{\partial\mu}
=\frac{e^2m^*}{\pi\hbar^2}.
\]
在理想二维抛物色散模型中这是常数，和 n 无关；若存在多子带、非抛物带、电子相互作用或有限温度，则实际量子电容可能随密度变化。
